\documentclass[10pt,a4paper]{amsart}
\usepackage{amsmath,amssymb,amsthm,graphicx}
% Retained preamble for the Stacks Project chapter morphisms.tex (Morphisms of Schemes).
% The chapter file pinned-original-source.tex carries only \input{preamble} on its first line, so no
% preamble text lives inside the chapter itself.  The Stacks Project uses one single shared preamble for
% every chapter; the snapshot stages it verbatim in sources/stacks-chow/preamble.tex.  The block below is
% copied mechanically, line for line, from that file's lines 14-18 (the Xy-pic setup) and lines 154-217
% (the native theorem declarations together with the mathematical macros).  Nothing was added, removed or
% redefined.  The wrapper class and the amsmath/amssymb/amsthm/graphicx base packages are supplied by the
% assembler, so the class/\externaldocument part of the shared preamble is not repeated here.
% For commutative diagrams we use Xy-pic
\usepackage[all]{xy}

% We use 2cell for 2-commutative diagrams.
\xyoption{2cell}
% Theorem environments.
%
\theoremstyle{plain}
\newtheorem{theorem}[subsection]{Theorem}
\newtheorem{proposition}[subsection]{Proposition}
\newtheorem{lemma}[subsection]{Lemma}

\theoremstyle{definition}
\newtheorem{definition}[subsection]{Definition}
\newtheorem{example}[subsection]{Example}
\newtheorem{exercise}[subsection]{Exercise}
\newtheorem{situation}[subsection]{Situation}

\theoremstyle{remark}
\newtheorem{remark}[subsection]{Remark}
\newtheorem{remarks}[subsection]{Remarks}

\numberwithin{equation}{subsection}

% Macros
%
\def\lim{\mathop{\mathrm{lim}}\nolimits}
\def\colim{\mathop{\mathrm{colim}}\nolimits}
\def\Spec{\mathop{\mathrm{Spec}}}
\def\Hom{\mathop{\mathrm{Hom}}\nolimits}
\def\Ext{\mathop{\mathrm{Ext}}\nolimits}
\def\SheafHom{\mathop{\mathcal{H}\!\mathit{om}}\nolimits}
\def\SheafExt{\mathop{\mathcal{E}\!\mathit{xt}}\nolimits}
\def\Sch{\mathit{Sch}}
\def\Mor{\mathop{\mathrm{Mor}}\nolimits}
\def\Ob{\mathop{\mathrm{Ob}}\nolimits}
\def\Sh{\mathop{\mathit{Sh}}\nolimits}
\def\NL{\mathop{N\!L}\nolimits}
\def\CH{\mathop{\mathrm{CH}}\nolimits}
\def\proetale{{pro\text{-}\acute{e}tale}}
\def\etale{{\acute{e}tale}}
\def\QCoh{\mathit{QCoh}}
\def\Ker{\mathop{\mathrm{Ker}}}
\def\Im{\mathop{\mathrm{Im}}}
\def\Coker{\mathop{\mathrm{Coker}}}
\def\Coim{\mathop{\mathrm{Coim}}}

% Boxtimes
%
\DeclareMathSymbol{\boxtimes}{\mathbin}{AMSa}{"02}

%
% Macros for moduli stacks/spaces
%
\def\QCohstack{\mathcal{QC}\!\mathit{oh}}
\def\Cohstack{\mathcal{C}\!\mathit{oh}}
\def\Spacesstack{\mathcal{S}\!\mathit{paces}}
\def\Quotfunctor{\mathrm{Quot}}
\def\Hilbfunctor{\mathrm{Hilb}}
\def\Curvesstack{\mathcal{C}\!\mathit{urves}}
\def\Polarizedstack{\mathcal{P}\!\mathit{olarized}}
\def\Complexesstack{\mathcal{C}\!\mathit{omplexes}}
% \Pic is the operator that assigns to X its picard group, usage \Pic(X)
% \Picardstack_{X/B} denotes the Picard stack of X over B
% \Picardfunctor_{X/B} denotes the Picard functor of X over B
\def\Pic{\mathop{\mathrm{Pic}}\nolimits}
\def\Picardstack{\mathcal{P}\!\mathit{ic}}
\def\Picardfunctor{\mathrm{Pic}}
\def\Deformationcategory{\mathcal{D}\!\mathit{ef}}

% --- appended 2026-10-12 (Chief): the project own definitions of the reference/slogan/history
% display environments, copied mechanically from the shared Stacks preamble lines 8-12
% (\usepackage{verbatim} plus the three \newenvironment lines, which the chapter preamble keeps
% verbatim from the project preamble). Nothing added, removed or redefined: these environments
% carry no mathematical content and the project itself typesets them as comments. Without these
% three lines a counted statement whose own text uses them fails with Environment slogan undefined.
\usepackage{verbatim}
\newenvironment{reference}{\comment}{\endcomment}
\newenvironment{slogan}{\comment}{\endcomment}
\newenvironment{history}{\comment}{\endcomment}
\usepackage{cleveref}
\crefname{originallemma}{original lemma}{original lemmas}
\Crefname{originallemma}{Original lemma}{Original lemmas}
\makeatletter
\newlabel{algebra-lemma-universal-omega}{{00RO}{source}{}{sourcepointer-1}{}}
\newlabel{algebra-lemma-universal-omega@cref}{{[originallemma][0][]00RO}{[1][0][]source}{}{}{}}
\makeatother
\crefname{originallemma}{original lemma}{original lemmas}
\Crefname{originallemma}{Original lemma}{Original lemmas}
\makeatletter
\newlabel{categories-lemma-yoneda}{{001P}{source}{}{sourcepointer-2}{}}
\newlabel{categories-lemma-yoneda@cref}{{[originallemma][0][]001P}{[1][0][]source}{}{}{}}
\makeatother
\crefname{originallemma}{original lemma}{original lemmas}
\Crefname{originallemma}{Original lemma}{Original lemmas}
\makeatletter
\newlabel{lemma-affine-case-derivation}{{01UO}{source}{}{sourcepointer-3}{}}
\newlabel{lemma-affine-case-derivation@cref}{{[originallemma][0][]01UO}{[1][0][]source}{}{}{}}
\makeatother
\crefname{originallemma}{original lemma}{original lemmas}
\Crefname{originallemma}{Original lemma}{Original lemmas}
\makeatletter
\newlabel{lemma-differentials-restrict-open}{{01US}{source}{}{sourcepointer-4}{}}
\newlabel{lemma-differentials-restrict-open@cref}{{[originallemma][0][]01US}{[1][0][]source}{}{}{}}
\makeatother
\crefname{originallemma}{original lemma}{original lemmas}
\Crefname{originallemma}{Original lemma}{Original lemmas}
\makeatletter
\newlabel{lemma-universal-derivation-universal}{{01UR}{source}{}{sourcepointer-5}{}}
\newlabel{lemma-universal-derivation-universal@cref}{{[originallemma][0][]01UR}{[1][0][]source}{}{}{}}
\makeatother
\crefname{originalsection}{original section}{original sections}
\Crefname{originalsection}{Original section}{Original sections}
\makeatletter
\newlabel{schemes-section-quasi-coherent-affine}{{01I6}{source}{}{sourcepointer-6}{}}
\newlabel{schemes-section-quasi-coherent-affine@cref}{{[originalsection][0][]01I6}{[1][0][]source}{}{}{}}
\makeatother
\title{The global sections of the sheaf of relative differentials over an affine chart are the module of differentials of the corresponding ring map}
\author{The Stacks Project Authors}
\date{stacks-project:morphisms.tex:local-snapshot-20261012; GNU FDL 1.2 or later}
\begin{document}
\maketitle
\noindent\textit{Editorial note.} All delivered mathematical text is the author's own, copied from the pinned Stacks Project chapter file pinned-original-source.tex; no statement or proof step was rewritten, reordered or removed.\par\medskip
\noindent\textit{Editorial note.} Following the multi-claim convention for this source, source\_url is the source-level URL https://stacks.math.columbia.edu/ and the source label lemma-differentials-affine is carried in claim\_key and here rather than in a per-claim tag URL.\par\medskip
\noindent\textit{Editorial note.} The preamble is the chapter's shared Stacks Project preamble, retained as sources/stacks-morphisms/retained-preamble.tex: a line-for-line excerpt of the shared Stacks preamble staged with this snapshot. No macro was invented or redefined, and this card adds no preamble text of its own.\par\medskip
\noindent\textit{Editorial note.} Statement span 6440-6449 and proof span 6451-6485 come from the halfA pool card operations/pool-cards/stacks-morphisms-halfA.json and were re-read in the pinned chapter file before this card was written.\par\medskip
\noindent\textit{Editorial note.} The displayed identifications of the three derivation and homomorphism groups are the author's own; they were copied unchanged.\par\medskip
\noindent\textit{Editorial note.} 6 cross-references occur inside the counted spans; they are listed in external\_references as navigation pointers only.\par\medskip
\noindent\textit{Allowed standard prerequisites (editorial).}\begin{itemize}
\item Restriction of the sheaf of relative differentials to open subschemes, cited from the same chapter. Cited at proof lines 6669-6670. Lets the proof replace X and S by the affine opens U and V and work with a ring map.
\item The universal property of the module of differentials and the affine description of derivations of a module sheaf. Cited at proof lines 6686-6688 and 6691-6692, the second being a card of this lane (L4005). Identifies the module at hand with the module of differentials of the ring map once it is shown to satisfy the universal property.
\item Yoneda's lemma and the universal property of the algebra module of differentials, cited from the categories and algebra chapters. Cited at proof lines 6698-6700. Converts the agreement of the represented functors into the asserted unique isomorphism.
\end{itemize}\medskip
\ifdefined\hypertarget\hypertarget{sourcepointer-1}{}\fi
\noindent\textit{Source reference.} Original lemma 00RO: focused-seeds/AGP/sources/stacks-snapshot-20261011/tags.txt, source lines 1010--1010. Same-chapter reference used in the retained author text at source line(s) 6481. It is supplied only as a navigation pointer so that the retained text compiles; the counted item delivers the author's own statement and the author's complete proof, and no part of the referenced item is copied in or substituted for the counted argument.\par\medskip
\ifdefined\hypertarget\hypertarget{sourcepointer-2}{}\fi
\noindent\textit{Source reference.} Original lemma 001P: focused-seeds/AGP/sources/stacks-snapshot-20261011/tags.txt, source lines 70--70. Same-chapter reference used in the retained author text at source line(s) 6482. It is supplied only as a navigation pointer so that the retained text compiles; the counted item delivers the author's own statement and the author's complete proof, and no part of the referenced item is copied in or substituted for the counted argument.\par\medskip
\ifdefined\hypertarget\hypertarget{sourcepointer-3}{}\fi
\noindent\textit{Source reference.} Original lemma 01UO: focused-seeds/AGP/sources/stacks-snapshot-20261011/tags.txt, source lines 2417--2417. Same-chapter reference used in the retained author text at source line(s) 6465. It is supplied only as a navigation pointer so that the retained text compiles; the counted item delivers the author's own statement and the author's complete proof, and no part of the referenced item is copied in or substituted for the counted argument.\par\medskip
\ifdefined\hypertarget\hypertarget{sourcepointer-4}{}\fi
\noindent\textit{Source reference.} Original lemma 01US: focused-seeds/AGP/sources/stacks-snapshot-20261011/tags.txt, source lines 2421--2421. Same-chapter reference used in the retained author text at source line(s) 6452. It is supplied only as a navigation pointer so that the retained text compiles; the counted item delivers the author's own statement and the author's complete proof, and no part of the referenced item is copied in or substituted for the counted argument.\par\medskip
\ifdefined\hypertarget\hypertarget{sourcepointer-5}{}\fi
\noindent\textit{Source reference.} Original lemma 01UR: focused-seeds/AGP/sources/stacks-snapshot-20261011/tags.txt, source lines 2420--2420. Same-chapter reference used in the retained author text at source line(s) 6464. It is supplied only as a navigation pointer so that the retained text compiles; the counted item delivers the author's own statement and the author's complete proof, and no part of the referenced item is copied in or substituted for the counted argument.\par\medskip
\ifdefined\hypertarget\hypertarget{sourcepointer-6}{}\fi
\noindent\textit{Source reference.} Original section 01I6: focused-seeds/AGP/sources/stacks-snapshot-20261011/tags.txt, source lines 1967--1967. Same-chapter reference used in the retained author text at source line(s) 6459. It is supplied only as a navigation pointer so that the retained text compiles; the counted item delivers the author's own statement and the author's complete proof, and no part of the referenced item is copied in or substituted for the counted argument.\par\medskip
\begin{lemma}
\label{lemma-differentials-affine}
Let $f : X \to S$ be a morphism of schemes. For any pair of affine opens
$\Spec(A) = U \subset X$, $\Spec(R) = V \subset S$ with $f(U) \subset V$
there is a unique isomorphism
$$
\Gamma(U, \Omega_{X/S}) = \Omega_{A/R}.
$$
compatible with $\text{d}_{X/S}$ and $\text{d} : A \to \Omega_{A/R}$.
\end{lemma}

\begin{proof}
By Lemma \ref{lemma-differentials-restrict-open} we may replace
$X$ and $S$ by $U$ and $V$. Thus we may assume $X = \Spec(A)$
and $S = \Spec(R)$ and we have to show the lemma with $U = X$
and $V = S$. Consider the $A$-module $M = \Gamma(X, \Omega_{X/S})$
together with the $R$-derivation $\text{d}_{X/S} : A \to M$.
Let $N$ be another $A$-module and denote $\widetilde{N}$ the quasi-coherent
$\mathcal{O}_X$-module associated to $N$, see
Schemes, Section \ref{schemes-section-quasi-coherent-affine}.
Precomposing by $\text{d}_{X/S} : A \to M$ we get an arrow
$$
\alpha : \Hom_A(M, N) \longrightarrow \text{Der}_R(A, N)
$$
Using Lemmas \ref{lemma-universal-derivation-universal} and
\ref{lemma-affine-case-derivation} we get identifications
$$
\Hom_{\mathcal{O}_X}(\Omega_{X/S}, \widetilde{N}) =
\text{Der}_S(\mathcal{O}_X, \widetilde{N}) =
\text{Der}_R(A, N)
$$
Taking global sections determines an arrow
$\Hom_{\mathcal{O}_X}(\Omega_{X/S}, \widetilde{N}) \to \Hom_R(M, N)$.
Combining this arrow and the identifications above we get an arrow
$$
\beta : \text{Der}_R(A, N) \longrightarrow \Hom_R(M, N)
$$
Checking what happens on global sections, we find that $\alpha$
and $\beta$ are each others inverse. Hence we see that
$\text{d}_{X/S} : A \to M$ satisfies
the same universal property as $\text{d} : A \to \Omega_{A/R}$,
see Algebra, Lemma \ref{algebra-lemma-universal-omega}. Thus the Yoneda lemma
(Categories, Lemma \ref{categories-lemma-yoneda})
implies there is a unique isomorphism of $A$-modules
$M \cong \Omega_{A/R}$ compatible with derivations.
\end{proof}

\end{document}
